\documentclass[11pt]{article}
\usepackage{amsmath, amssymb, geometry}
\geometry{a4paper, margin=1in}

\begin{document}

\title{Homework 15.1: Double Integrals over Rectangles - Practice Set}
\author{Math Practice Generator}
\date{}
\maketitle

\section*{Problem Set: Double Integrals over Rectangles}

\noindent\textbf{Problem 1: Geometric Volume of a Constant Height} \\
\textbf{Question:} Evaluate the double integral by identifying it as the volume of a solid: $\iint_R 4 \, dA$, where $R = \{(x,y) \mid 2 \leq x \leq 5, 0 \leq y \leq 3\}$. \\
\textbf{Solution:} This represents the volume of a rectangular box with a constant height $z = 4$ over the region $R$. 
1. Dimensions of the base $R$: Length = $5 - 2 = 3$. Width = $3 - 0 = 3$.
2. Area of base $A = 3 \times 3 = 9$.
3. Volume = Area $\times$ Height = $9 \times 4 = 36$. \\
\textbf{Answer:} 36

\vspace{0.5cm}
\noindent\textbf{Problem 2: Geometric Volume of a Plane} \\
\textbf{Question:} Evaluate the double integral $\iint_R \pi \, dA$, $R = [0,1] \times [0,1]$ using geometry. \\
\textbf{Solution:} The function is a constant $z = \pi$. The base is a $1 \times 1$ square.
1. Area of base $A = 1 \times 1 = 1$.
2. Volume = $1 \times \pi = \pi$. \\
\textbf{Answer:} $\pi$

\vspace{0.5cm}
\noindent\textbf{Problem 3: Riemann Sum (Upper Right Corners)} \\
\textbf{Question:} Estimate the volume under $z = xy^2$ over $R = [0,4] \times [0,2]$ using a Riemann Sum with $m=2, n=2$ and upper right corner sample points. \\
\textbf{Solution:} 
$\Delta x = \frac{4-0}{2} = 2$, $\Delta y = \frac{2-0}{2} = 1$. Area $\Delta A = 2 \times 1 = 2$.
Grid intervals: $x \in [0,2], [2,4]$; $y \in [0,1], [1,2]$. 
Upper Right points $(x_i, y_j)$: $(2,1), (4,1), (2,2), (4,2)$.
Evaluations: $f(2,1) = 2(1)^2 = 2$; $f(4,1) = 4(1)^2 = 4$; $f(2,2) = 2(2)^2 = 8$; $f(4,2) = 4(2)^2 = 16$.
Sum of values = $2 + 4 + 8 + 16 = 30$.
Volume $\approx 30 \times \Delta A = 30 \times 2 = 60$. \\
\textbf{Answer:} 60

\vspace{0.5cm}
\noindent\textbf{Problem 4: Riemann Sum (Lower Left Corners)} \\
\textbf{Question:} Estimate the volume under $z = x^2+y$ over $R = [1,3] \times [1,5]$ with $m=2, n=2$, using lower left corner sample points. \\
\textbf{Solution:} 
$\Delta x = \frac{3-1}{2} = 1$, $\Delta y = \frac{5-1}{2} = 2$. Area $\Delta A = 1 \times 2 = 2$.
Grid intervals: $x \in [1,2], [2,3]$; $y \in [1,3], [3,5]$.
Lower Left points: $(1,1), (2,1), (1,3), (2,3)$.
Evaluations: $f(1,1) = 1+1=2$; $f(2,1) = 4+1=5$; $f(1,3) = 1+3=4$; $f(2,3) = 4+3=7$.
Sum = $2 + 5 + 4 + 7 = 18$.
Volume $\approx 18 \times \Delta A = 18 \times 2 = 36$. \\
\textbf{Answer:} 36

\vspace{0.5cm}
\noindent\textbf{Problem 5: Riemann Sum (Midpoint Rule)} \\
\textbf{Question:} Estimate the integral of $f(x,y) = x-y$ over $R = [0,4] \times [0,4]$ using the Midpoint Rule with $m=2, n=2$. \\
\textbf{Solution:} 
$\Delta x = 2$, $\Delta y = 2$, $\Delta A = 4$.
Sub-rectangles are centered at midpoints: $(1,1), (3,1), (1,3), (3,3)$.
Evaluations: $f(1,1) = 1-1=0$; $f(3,1) = 3-1=2$; $f(1,3) = 1-3=-2$; $f(3,3) = 3-3=0$.
Sum = $0 + 2 - 2 + 0 = 0$.
Volume $\approx 0 \times 4 = 0$. \\
\textbf{Answer:} 0

\vspace{0.5cm}
\noindent\textbf{Problem 6: Exponential Riemann Sum} \\
\textbf{Question:} Estimate $\iint_R e^{x+y} dA$ over $R = [0,2] \times [0,1]$ using a Riemann sum with $m=2, n=1$ and upper right corners. \\
\textbf{Solution:} 
$\Delta x = \frac{2-0}{2} = 1$, $\Delta y = \frac{1-0}{1} = 1$, $\Delta A = 1$.
Upper right sample points: $(1,1), (2,1)$.
Evaluations: $f(1,1) = e^{1+1} = e^2$; $f(2,1) = e^{2+1} = e^3$.
Sum $\approx (e^2 + e^3) \times 1 = e^2 + e^3$. \\
\textbf{Answer:} $e^2 + e^3$

\vspace{0.5cm}
\noindent\textbf{Problem 7: Partial Integral with respect to y} \\
\textbf{Question:} Evaluate the partial integral $\int_1^3 (4x^3 y - 2e^x) \, dy$. \\
\textbf{Solution:} Treat $x$ as a constant.
$\int_1^3 (4x^3 y - 2e^x) dy = \left[ 2x^3 y^2 - 2y e^x \right]_{y=1}^{y=3}$
Upper bound: $2x^3(9) - 2(3)e^x = 18x^3 - 6e^x$.
Lower bound: $2x^3(1) - 2(1)e^x = 2x^3 - 2e^x$.
Result: $(18x^3 - 6e^x) - (2x^3 - 2e^x) = 16x^3 - 4e^x$. \\
\textbf{Answer:} $16x^3 - 4e^x$

\vspace{0.5cm}
\noindent\textbf{Problem 8: Partial Integral with respect to x} \\
\textbf{Question:} Evaluate the partial integral $\int_0^{\pi} (y \sin x + y^2) \, dx$. \\
\textbf{Solution:} Treat $y$ as a constant.
$\int_0^{\pi} (y \sin x + y^2) dx = \left[ -y \cos x + xy^2 \right]_{0}^{\pi}$
Upper bound: $-y \cos(\pi) + \pi y^2 = y + \pi y^2$.
Lower bound: $-y \cos(0) + 0 = -y$.
Result: $(y + \pi y^2) - (-y) = 2y + \pi y^2$. \\
\textbf{Answer:} $2y + \pi y^2$

\vspace{0.5cm}
\noindent\textbf{Problem 9: Iterated Integral Setup and Evaluation} \\
\textbf{Question:} Evaluate $\int_0^2 \int_0^1 12x^2 y^3 \, dx \, dy$. \\
\textbf{Solution:} 
Inner (wrt $x$): $\int_0^1 12x^2 y^3 dx = \left[ 4x^3 y^3 \right]_0^1 = 4y^3$.
Outer (wrt $y$): $\int_0^2 4y^3 dy = \left[ y^4 \right]_0^2 = 16 - 0 = 16$. \\
\textbf{Answer:} 16

\vspace{0.5cm}
\noindent\textbf{Problem 10: Iterated Integral with Polynomials} \\
\textbf{Question:} Evaluate $\int_1^2 \int_0^3 (2x - 4y) \, dy \, dx$. \\
\textbf{Solution:} 
Inner (wrt $y$): $\int_0^3 (2x - 4y) dy = \left[ 2xy - 2y^2 \right]_0^3 = (6x - 18) - 0$.
Outer (wrt $x$): $\int_1^2 (6x - 18) dx = \left[ 3x^2 - 18x \right]_1^2 = (12 - 36) - (3 - 18) = -24 - (-15) = -9$. \\
\textbf{Answer:} -9

\vspace{0.5cm}
\noindent\textbf{Problem 11: Iterated Order $dy \, dx$} \\
\textbf{Question:} Evaluate $\int_{-1}^1 \int_{-2}^2 (x^2 + y^2) \, dy \, dx$. \\
\textbf{Solution:} 
Inner (wrt $y$): $\int_{-2}^2 (x^2 + y^2) dy = \left[ x^2 y + \frac{y^3}{3} \right]_{-2}^2 = (2x^2 + \frac{8}{3}) - (-2x^2 - \frac{8}{3}) = 4x^2 + \frac{16}{3}$.
Outer (wrt $x$): $\int_{-1}^1 (4x^2 + \frac{16}{3}) dx = \left[ \frac{4x^3}{3} + \frac{16x}{3} \right]_{-1}^1 = (\frac{4}{3} + \frac{16}{3}) - (-\frac{4}{3} - \frac{16}{3}) = 2(\frac{20}{3}) = \frac{40}{3}$. \\
\textbf{Answer:} 40/3

\vspace{0.5cm}
\noindent\textbf{Problem 12: Separable Trigonometric Integral} \\
\textbf{Question:} Evaluate $\int_0^{\pi/2} \int_0^{\pi/2} \cos(x)\sin(y) \, dy \, dx$. \\
\textbf{Solution:} Since the function is separable and bounds are constants, write as a product:
$\left(\int_0^{\pi/2} \cos x \, dx\right) \left(\int_0^{\pi/2} \sin y \, dy\right)$.
1. $\int_0^{\pi/2} \cos x \, dx = [\sin x]_0^{\pi/2} = 1 - 0 = 1$.
2. $\int_0^{\pi/2} \sin y \, dy = [-\cos y]_0^{\pi/2} = 0 - (-1) = 1$.
Total: $1 \times 1 = 1$. \\
\textbf{Answer:} 1

\vspace{0.5cm}
\noindent\textbf{Problem 13: Separable Function with Negative Exponent} \\
\textbf{Question:} Evaluate $\int_1^3 \int_1^2 \frac{x^2}{y^3} \, dy \, dx$. \\
\textbf{Solution:} Separate variables: $\left(\int_1^3 x^2 \, dx\right) \left(\int_1^2 y^{-3} \, dy\right)$.
1. $\left[\frac{x^3}{3}\right]_1^3 = 9 - \frac{1}{3} = \frac{26}{3}$.
2. $\left[\frac{y^{-2}}{-2}\right]_1^2 = -\frac{1}{2}(\frac{1}{4} - 1) = -\frac{1}{2}(-\frac{3}{4}) = \frac{3}{8}$.
Total: $\frac{26}{3} \times \frac{3}{8} = \frac{78}{24} = \frac{13}{4}$. \\
\textbf{Answer:} 13/4

\vspace{0.5cm}
\noindent\textbf{Problem 14: Separable Exponential Integral} \\
\textbf{Question:} Evaluate $\iint_R e^{2x+y} \, dA$, where $R = [0,\ln 2] \times [0,\ln 3]$. \\
\textbf{Solution:} Write as $e^{2x}e^y$ and separate:
$\left(\int_0^{\ln 2} e^{2x} dx\right)\left(\int_0^{\ln 3} e^y dy\right)$.
1. $\left[\frac{1}{2}e^{2x}\right]_0^{\ln 2} = \frac{1}{2}(e^{2\ln 2} - 1) = \frac{1}{2}(4 - 1) = \frac{3}{2}$.
2. $\left[e^y\right]_0^{\ln 3} = e^{\ln 3} - 1 = 3 - 1 = 2$.
Total: $\frac{3}{2} \times 2 = 3$. \\
\textbf{Answer:} 3

\vspace{0.5cm}
\noindent\textbf{Problem 15: Separable Fractional Powers} \\
\textbf{Question:} Evaluate $\iint_R x \sqrt{y} \, dA$, where $R = [0,2] \times [1,4]$. \\
\textbf{Solution:} Separate variables: $\left(\int_0^2 x \, dx\right)\left(\int_1^4 y^{1/2} \, dy\right)$.
1. $\left[\frac{x^2}{2}\right]_0^2 = \frac{4}{2} = 2$.
2. $\left[\frac{2}{3}y^{3/2}\right]_1^4 = \frac{2}{3}(4^{3/2} - 1^{3/2}) = \frac{2}{3}(8 - 1) = \frac{14}{3}$.
Total: $2 \times \frac{14}{3} = \frac{28}{3}$. \\
\textbf{Answer:} 28/3

\vspace{0.5cm}
\noindent\textbf{Problem 16: U-Substitution (Inner Integral)} \\
\textbf{Question:} Evaluate $\int_0^1 \int_0^{\pi} x \cos(xy) \, dy \, dx$. \\
\textbf{Solution:} Inner (wrt $y$): Let $u = xy \implies du = x \, dy$. The integral of $x \cos(xy)$ wrt $y$ is $\sin(xy)$.
Evaluate from $y=0$ to $y=\pi$: $\left[ \sin(xy) \right]_0^\pi = \sin(\pi x) - 0 = \sin(\pi x)$.
Outer (wrt $x$): $\int_0^1 \sin(\pi x) \, dx = \left[ -\frac{1}{\pi} \cos(\pi x) \right]_0^1 = -\frac{1}{\pi}(\cos(\pi) - \cos(0)) = -\frac{1}{\pi}(-1 - 1) = \frac{2}{\pi}$. \\
\textbf{Answer:} $2/\pi$

\vspace{0.5cm}
\noindent\textbf{Problem 17: U-Substitution (Iterated Evaluation)} \\
\textbf{Question:} Evaluate $\int_0^1 \int_0^1 y e^{xy} \, dx \, dy$. \\
\textbf{Solution:} Inner (wrt $x$): Let $u = xy \implies du = y \, dx$. The integral of $y e^{xy}$ wrt $x$ is $e^{xy}$.
$\left[ e^{xy} \right]_0^1 = e^{y(1)} - e^0 = e^y - 1$.
Outer (wrt $y$): $\int_0^1 (e^y - 1) \, dy = \left[ e^y - y \right]_0^1 = (e^1 - 1) - (1 - 0) = e - 2$. \\
\textbf{Answer:} $e - 2$

\vspace{0.5cm}
\noindent\textbf{Problem 18: Algebraic U-Substitution} \\
\textbf{Question:} Evaluate $\int_0^2 \int_0^1 2x \sqrt{x^2+y} \, dx \, dy$. \\
\textbf{Solution:} Inner (wrt $x$): Let $u = x^2+y \implies du = 2x \, dx$. Integral is $\int u^{1/2} du = \frac{2}{3}u^{3/2}$.
$\left[ \frac{2}{3}(x^2+y)^{3/2} \right]_0^1 = \frac{2}{3}((1+y)^{3/2} - y^{3/2})$.
Outer (wrt $y$): $\frac{2}{3} \int_0^2 ((1+y)^{3/2} - y^{3/2}) \, dy$.
$\frac{2}{3} \left[ \frac{2}{5}(1+y)^{5/2} - \frac{2}{5}y^{5/2} \right]_0^2 = \frac{4}{15} \left( (3^{5/2} - 2^{5/2}) - (1^{5/2} - 0) \right) = \frac{4}{15}(9\sqrt{3} - 4\sqrt{2} - 1)$. \\
\textbf{Answer:} $\frac{4}{15}(9\sqrt{3} - 4\sqrt{2} - 1)$

\vspace{0.5cm}
\noindent\textbf{Problem 19: Logarithmic Integration} \\
\textbf{Question:} Evaluate $\int_0^1 \int_0^1 \frac{x}{1+xy} \, dy \, dx$. \\
\textbf{Solution:} Inner (wrt $y$): Let $u = 1+xy \implies du = x \, dy$. Integral of $\frac{1}{u} du$ is $\ln|u|$.
$\left[ \ln|1+xy| \right]_0^1 = \ln(1+x) - \ln(1) = \ln(1+x)$.
Outer (wrt $x$): $\int_0^1 \ln(1+x) \, dx$. Use integration by parts. Let $w = 1+x, dw = dx$.
$\int_1^2 \ln(w) dw = [w \ln w - w]_1^2 = (2 \ln 2 - 2) - (1 \ln 1 - 1) = 2 \ln 2 - 1$. \\
\textbf{Answer:} $2 \ln 2 - 1$

\vspace{0.5cm}
\noindent\textbf{Problem 20: Negative Logarithmic Bounds} \\
\textbf{Question:} Evaluate $\iint_R \frac{2}{xy^3} \, dA$, where $R = [-e,-1] \times [1,2]$. \\
\textbf{Solution:} Separable: $\left(\int_{-e}^{-1} \frac{2}{x} dx\right)\left(\int_1^2 y^{-3} dy\right)$.
1. $\left[2\ln|x|\right]_{-e}^{-1} = 2(\ln|-1| - \ln|-e|) = 2(0 - 1) = -2$. (Crucial: absolute value of negative bounds).
2. $\left[\frac{y^{-2}}{-2}\right]_1^2 = -\frac{1}{2}(\frac{1}{4} - 1) = \frac{3}{8}$.
Total: $-2 \times \frac{3}{8} = -\frac{3}{4}$. \\
\textbf{Answer:} -3/4

\vspace{0.5cm}
\noindent\textbf{Problem 21: Rational/Arctan Integration} \\
\textbf{Question:} Evaluate $\iint_R \frac{y^2}{x^2+1} \, dA$, $R = [0,1] \times [0,3]$. \\
\textbf{Solution:} Separable: $\left(\int_0^1 \frac{1}{x^2+1} dx\right)\left(\int_0^3 y^2 dy\right)$.
1. $\left[\arctan x\right]_0^1 = \frac{\pi}{4} - 0 = \frac{\pi}{4}$.
2. $\left[\frac{y^3}{3}\right]_0^3 = \frac{27}{3} = 9$.
Total: $\frac{9\pi}{4}$. \\
\textbf{Answer:} $\frac{9\pi}{4}$

\vspace{0.5cm}
\noindent\textbf{Problem 22: Integration by Parts (Trigonometric)} \\
\textbf{Question:} Evaluate $\iint_R 2x \sin(x-y) \, dA$, $R = [0,\pi/2] \times [0,\pi/2]$. \\
\textbf{Solution:} Inner (wrt $y$): $\int_0^{\pi/2} 2x \sin(x-y) \, dy = \left[ 2x \cos(x-y) \right]_0^{\pi/2} = 2x \cos(x-\pi/2) - 2x \cos(x)$.
Using $\cos(x-\pi/2) = \sin(x)$, this simplifies to $2x \sin x - 2x \cos x$.
Outer (wrt $x$): $\int_0^{\pi/2} (2x \sin x - 2x \cos x) \, dx$.
Use IBP on both terms. $\int x \sin x = -x \cos x + \sin x$ and $\int x \cos x = x \sin x + \cos x$.
Resulting evaluated expression: $\left[ -2x\cos x + 2\sin x - 2x\sin x - 2\cos x \right]_0^{\pi/2}$.
At $\pi/2$: $0 + 2 - \pi - 0 = 2 - \pi$.
At $0$: $0 + 0 - 0 - 2 = -2$.
Difference: $(2-\pi) - (-2) = 4 - \pi$. \\
\textbf{Answer:} $4 - \pi$

\vspace{0.5cm}
\noindent\textbf{Problem 23: Log Integration and IBP} \\
\textbf{Question:} Evaluate $\iint_R \frac{3y}{1+xy} \, dA$, $R = [0,1] \times [0,2]$. \\
\textbf{Solution:} Inner (wrt $x$): Let $u = 1+xy, du = y dx$. Integral is $3\ln(1+xy)$.
Evaluate from 0 to 1: $3\ln(1+y) - 3\ln(1) = 3\ln(1+y)$.
Outer (wrt $y$): $\int_0^2 3\ln(1+y) \, dy$. Let $w = 1+y \implies dw = dy$. $\int_1^3 3\ln(w) dw$.
IBP: $3[w \ln w - w]_1^3 = 3(3\ln 3 - 3) - 3(0 - 1) = 9\ln 3 - 9 + 3 = 9\ln 3 - 6$. \\
\textbf{Answer:} $9\ln 3 - 6$

\vspace{0.5cm}
\noindent\textbf{Problem 24: Symmetry Shortcut 1} \\
\textbf{Question:} Evaluate $\iint_R x^3 \cos(y) \, dA$, $R = [-2,2] \times [0,\pi]$. \\
\textbf{Solution:} The function is separable into $f(x) = x^3$ and $g(y) = \cos(y)$.
Since $f(x) = x^3$ is an odd function ($f(-x) = -f(x)$) being integrated over a symmetric interval $[-2,2]$, $\int_{-2}^2 x^3 dx = 0$.
Total integral = $0 \times \int_0^\pi \cos y dy = 0$. \\
\textbf{Answer:} 0

\vspace{0.5cm}
\noindent\textbf{Problem 25: Symmetry Shortcut 2} \\
\textbf{Question:} Evaluate $\iint_R \frac{\sin(x)}{y^2+1} \, dA$, $R = [-\pi/4, \pi/4] \times [0,1]$. \\
\textbf{Solution:} Separable: $\left(\int_{-\pi/4}^{\pi/4} \sin(x) dx\right)\left(\int_0^1 \frac{1}{y^2+1} dy\right)$.
The sine function is odd, and the integration interval is symmetric around zero. Thus, the integral with respect to $x$ evaluates to 0. 
Total integral = 0. \\
\textbf{Answer:} 0

\vspace{0.5cm}
\noindent\textbf{Problem 26: Choosing the Order of Integration} \\
\textbf{Question:} Evaluate $\iint_R y \cos(xy) \, dA$, $R = [0,1] \times [0,\pi/2]$. Set up the order $dx \, dy$ to avoid integration by parts. \\
\textbf{Solution:} If we integrate wrt $x$ first, we can use simple u-substitution.
Inner (wrt $x$): $\int_0^1 y \cos(xy) \, dx = \left[ \sin(xy) \right]_0^1 = \sin(y) - \sin(0) = \sin(y)$.
Outer (wrt $y$): $\int_0^{\pi/2} \sin(y) \, dy = \left[ -\cos(y) \right]_0^{\pi/2} = 0 - (-1) = 1$. \\
\textbf{Answer:} 1

\vspace{0.5cm}
\noindent\textbf{Problem 27: Average Value of a Function 1} \\
\textbf{Question:} Find the average value of $f(x,y) = x^2 y$ over the rectangle $R = [0,3] \times [0,4]$. \\
\textbf{Solution:} Formula for average value: $f_{ave} = \frac{1}{\text{Area}(R)} \iint_R f(x,y) dA$.
Area $= 3 \times 4 = 12$.
Integral (Separable): $\left(\int_0^3 x^2 dx\right)\left(\int_0^4 y dy\right) = \left[\frac{x^3}{3}\right]_0^3 \left[\frac{y^2}{2}\right]_0^4 = (9)(8) = 72$.
Average Value $= \frac{72}{12} = 6$. \\
\textbf{Answer:} 6

\vspace{0.5cm}
\noindent\textbf{Problem 28: Average Value of a Function 2} \\
\textbf{Question:} Find the average value of $f(x,y) = e^y \sqrt{x}$ over $R = [0,4] \times [0,\ln 5]$. \\
\textbf{Solution:} Area $= 4 \times \ln 5 = 4\ln 5$.
Integral (Separable): $\left(\int_0^4 x^{1/2} dx\right)\left(\int_0^{\ln 5} e^y dy\right) = \left[\frac{2}{3}x^{3/2}\right]_0^4 \left[e^y\right]_0^{\ln 5}$.
$= \frac{2}{3}(8) \times (5 - 1) = \frac{16}{3} \times 4 = \frac{64}{3}$.
Average Value $= \frac{64/3}{4\ln 5} = \frac{16}{3\ln 5}$. \\
\textbf{Answer:} $\frac{16}{3\ln 5}$

\vspace{0.5cm}
\noindent\textbf{Problem 29: Advanced Diagnostic: Find the Flaw} \\
\textbf{Question:} Evaluate $\iint_R (x^2+y^2) \, dA$, $R = [0,1] \times [0,1]$. A flawed student separates this as $\left(\int_0^1 x^2 dx\right)\left(\int_0^1 y^2 dy\right) = 1/9$. Find the correct answer and identify the flaw. \\
\textbf{Solution:} The flaw is assuming $x^2+y^2$ is separable. Separability only works for products $g(x)h(y)$, not sums.
Correct approach is an iterated integral:
Inner (wrt $y$): $\int_0^1 (x^2 + y^2) dy = \left[ x^2 y + \frac{y^3}{3} \right]_0^1 = x^2 + \frac{1}{3}$.
Outer (wrt $x$): $\int_0^1 (x^2 + \frac{1}{3}) dx = \left[ \frac{x^3}{3} + \frac{x}{3} \right]_0^1 = \frac{1}{3} + \frac{1}{3} = \frac{2}{3}$. \\
\textbf{Answer:} $2/3$

\vspace{0.5cm}
\noindent\textbf{Problem 30: Mixed Techniques with Negative Exponents} \\
\textbf{Question:} Evaluate $\int_1^2 \int_1^2 \left( \frac{x^2}{y^2} + y e^x \right) \, dy \, dx$. \\
\textbf{Solution:} Note that $\frac{1}{y^2} = y^{-2}$. 
Inner (wrt $y$): $\int_1^2 (x^2 y^{-2} + y e^x) dy = \left[ -x^2 y^{-1} + \frac{y^2}{2} e^x \right]_1^2$
$= \left( -\frac{x^2}{2} + 2e^x \right) - \left( -x^2 + \frac{1}{2}e^x \right) = \frac{x^2}{2} + \frac{3}{2}e^x$.
Outer (wrt $x$): $\int_1^2 \left( \frac{x^2}{2} + \frac{3}{2}e^x \right) dx = \left[ \frac{x^3}{6} + \frac{3}{2}e^x \right]_1^2$
$= \left( \frac{8}{6} + \frac{3}{2}e^2 \right) - \left( \frac{1}{6} + \frac{3}{2}e \right) = \frac{7}{6} + \frac{3}{2}e^2 - \frac{3}{2}e$. \\
\textbf{Answer:} $\frac{7}{6} + \frac{3}{2}(e^2 - e)$

\newpage
\section*{Concept Checklist}
The following outlines the concepts covered across the practice set matching the curriculum described in the source document:

\begin{itemize}
    \item \textbf{Geometric Volume Interpretation} (Volume of solid, plane): Problems 1, 2
    \item \textbf{Riemann Sums}
    \begin{itemize}
        \item Upper Right Corners: Problems 3, 6
        \item Lower Left Corners: Problem 4
        \item Midpoint Rule: Problem 5
    \end{itemize}
    \item \textbf{Partial Integration} (Treating $x$ or $y$ as constants): Problems 7, 8
    \item \textbf{Iterated Integrals} (Fubini's Theorem, basic polynomials): Problems 9, 10, 11
    \item \textbf{Separable Functions} ($f(x,y) = g(x)h(y)$): Problems 12, 13, 14, 15, 20, 21
    \item \textbf{U-Substitution Methods}: Problems 16, 17, 18, 19
    \item \textbf{Integration By Parts}: Problems 19, 22, 23
    \item \textbf{Logarithmic Integration \& Absolute Value Bounds}: Problems 19, 20, 23
    \item \textbf{Symmetry Tricks} (Odd function over $[-a, a]$ bounds $= 0$): Problems 24, 25
    \item \textbf{Order of Integration Optimization} ($dx \, dy$ vs $dy \, dx$ to simplify work): Problem 26
    \item \textbf{Average Value of a Function} ($f_{ave} = \frac{1}{A} \iint f \, dA$): Problems 27, 28
    \item \textbf{Conceptual Flaws / Untested Fallacies} (Separability abuse traps): Problem 29
    \item \textbf{Negative Exponents and Mixed Algebraic Expressions}: Problems 13, 20, 30
\end{itemize}

\end{document}